\documentclass[11pt,twocolumn]{article}

\usepackage[margin=1in]{geometry}
\usepackage{amsmath,amssymb,amsthm}
\usepackage{graphicx}
\usepackage{hyperref}
\usepackage{listings}
\usepackage{xcolor}
\usepackage{booktabs}
\usepackage{enumitem}

\hypersetup{colorlinks=true,linkcolor=blue!60!black,citecolor=green!40!black,urlcolor=blue!60!black}

\lstset{
  language=Rust,
  basicstyle=\ttfamily\scriptsize,
  keywordstyle=\color{blue!70!black}\bfseries,
  commentstyle=\color{green!50!black}\itshape,
  stringstyle=\color{red!60!black},
  breaklines=true,
  frame=single,
  backgroundcolor=\color{gray!5}
}

\newtheorem{definition}{Definition}
\newtheorem{theorem}{Theorem}
\newtheorem{proposition}{Proposition}

\title{Decentralized Knowledge Verification: Claims, Inference, and Factchecking with Epistemic Mesh}
\author{Tristan Stoltz \\ Luminous Dynamics \\ \texttt{tristan.stoltz@evolvingresonantcocreationism.com}}
\date{March 2026}

\begin{document}

\twocolumn[
\begin{@twocolumnfalse}
\maketitle

\begin{abstract}
We present a decentralized knowledge verification system implemented as 8
Holochain zomes: claims registration, knowledge graph construction, semantic
query, logical inference, multi-source factchecking, prediction market
integration, distributed key generation for knowledge attestation, and
cross-cluster bridge. The system forms an ``epistemic mesh'' where knowledge
claims are verified through multiple independent pathways---peer review,
logical inference, prediction market pricing, and cross-domain
corroboration---each contributing to a composite confidence score. Verification
quality is weighted by the verifier's consciousness tier and domain engagement,
creating an incentive structure where epistemic authority derives from
demonstrated knowledge rather than institutional position. The knowledge graph
is distributed across Holochain's DHT with typed relationships (supports,
contradicts, extends, specializes) and supports transitive confidence
propagation across claim chains. Integration with 15 other Mycelix clusters
enables cross-domain knowledge verification: a climate claim can be corroborated
by supply chain provenance data, space-based observations, and energy grid
measurements.
\end{abstract}
\end{@twocolumnfalse}
]

\section{Introduction}

Centralized knowledge systems---Wikipedia, Snopes, institutional peer
review---face a fundamental trust problem: who verifies the verifiers?
Centralized platforms concentrate epistemic authority in editors,
administrators, and reviewers whose incentives may not align with truth.

Decentralized alternatives have emerged---prediction markets for forecasting,
blockchain-based fact-checking protocols, and peer-to-peer review
systems---but these typically address only one aspect of knowledge
verification. A claim about climate change requires scientific evidence,
statistical analysis, observational data, and economic context. No single
verification mechanism suffices.

We present the Mycelix Knowledge cluster: an \emph{epistemic mesh} where
multiple verification pathways converge on composite confidence scores for
knowledge claims.

\section{Architecture}

\subsection{Seven Zomes plus Bridge}

\begin{enumerate}[nosep]
  \item \textbf{claims}: Structured claim registration with provenance
    metadata, evidence links, domain tags, and confidence scoring.
  \item \textbf{graph}: Knowledge graph construction with typed
    relationships between claims.
  \item \textbf{query}: Semantic query processing across the distributed
    knowledge graph.
  \item \textbf{inference}: Logical inference engine for deriving new
    claims from existing premises.
  \item \textbf{factcheck}: Multi-source fact verification with
    consciousness-weighted confidence aggregation.
  \item \textbf{markets\_integration}: Prediction market integration for
    epistemic value discovery.
  \item \textbf{bridge}: Cross-cluster knowledge dispatch.
\end{enumerate}

\subsection{Data Model}

Claims are the atomic unit of knowledge:

\begin{lstlisting}
pub struct Claim {
    pub statement: String,
    pub domain: Vec<String>,
    pub evidence: Vec<EvidenceLink>,
    pub confidence: f64,  // [0, 1]
    pub author: AgentPubKey,
    pub created_at: Timestamp,
    pub schema_version: u8,
}

pub struct EvidenceLink {
    pub source: EvidenceSource,
    pub uri: String,
    pub strength: f64,  // [0, 1]
}

pub enum EvidenceSource {
    DirectObservation,
    PeerReview,
    StatisticalAnalysis,
    ExpertTestimony,
    InstrumentMeasurement,
    CrossDomainCorroboration,
}
\end{lstlisting}

\section{Knowledge Graph}

\subsection{Typed Relationships}

The graph zome maintains a typed knowledge graph where nodes are claims
and edges are relationships:

\begin{definition}[Knowledge Edge Types]
\begin{enumerate}[nosep]
  \item \textbf{Supports}: Claim A provides evidence for Claim B.
    Strength: how strongly A supports B.
  \item \textbf{Contradicts}: Claim A provides counter-evidence to Claim B.
    Strength: how strongly A undermines B.
  \item \textbf{Extends}: Claim A generalizes or broadens Claim B.
  \item \textbf{Specializes}: Claim A is a specific instance of Claim B.
  \item \textbf{Requires}: Claim A is a necessary premise for Claim B.
  \item \textbf{Supersedes}: Claim A replaces an outdated Claim B.
\end{enumerate}
\end{definition}

\subsection{DHT Distribution}

The knowledge graph is distributed across Holochain's DHT:

\begin{itemize}[nosep]
  \item Claim entries are hashed and stored in their DHT neighborhood
  \item Link entries (relationships) are attached to source claim hashes
  \item Query processing traverses links across DHT neighborhoods
  \item Validation rules ensure relationship consistency (e.g., a
    Supersedes edge requires both claims to share at least one domain tag)
\end{itemize}

\subsection{Transitive Confidence Propagation}

Confidence propagates through the knowledge graph along Support edges:

\begin{equation}
  c_{\text{derived}}(B) = \max\left(c_{\text{direct}}(B), \;
    \max_{A \to_s B} c(A) \cdot s(A, B) \cdot d^{h(A,B)}\right)
\end{equation}

where $s(A, B)$ is the support strength, $d \in (0, 1)$ is the per-hop
decay, and $h(A, B)$ is the hop distance. Contradiction edges reduce
confidence:

\begin{equation}
  c_{\text{adjusted}}(B) = c_{\text{derived}}(B) \cdot
    \prod_{A \to_c B} (1 - c(A) \cdot s_c(A, B))
\end{equation}

This formulation ensures that a single strong contradiction can significantly
reduce confidence, while weak contradictions have proportionally small effects.

\section{Claim Lifecycle}

\subsection{Registration}

Any Participant+ tier agent can register a claim. The claim entry is
validated by the integrity zome:

\begin{enumerate}[nosep]
  \item Statement is non-empty and within length limits
  \item At least one domain tag is provided
  \item Author's consciousness credential is valid
  \item Schema version is current
\end{enumerate}

\subsection{Verification}

Independent agents verify claims by submitting verification entries:

\begin{lstlisting}
pub struct Verification {
    pub claim_hash: ActionHash,
    pub verdict: Verdict,
    pub evidence: Vec<EvidenceLink>,
    pub confidence: f64,
    pub verifier: AgentPubKey,
}

pub enum Verdict {
    Confirmed,
    Disputed,
    Inconclusive,
    NeedsMoreEvidence,
}
\end{lstlisting}

Verification quality is weighted by the verifier's consciousness profile:

\begin{equation}
  w_j = E_j \cdot \text{tier\_weight}(\text{tier}_j)
\end{equation}

where $E_j$ is the verifier's engagement in the claim's domain and
$\text{tier\_weight}$ maps consciousness tiers to weights. This ensures
that domain experts with strong governance standing contribute more to
the consensus.

\subsection{Inference}

The inference zome derives new claims from existing premises using
logical rules:

\begin{lstlisting}
pub struct InferenceRule {
    pub premises: Vec<ClaimPattern>,
    pub conclusion: ClaimTemplate,
    pub confidence_function: ConfidenceFn,
}
\end{lstlisting}

Inference runs periodically (not in real-time), scanning the knowledge
graph for premise patterns and generating derived claims with
appropriate provenance metadata. Derived claims are marked as
inference-generated and link back to their premises via Requires edges.

\subsection{Factchecking}

The factcheck zome aggregates verification results:

\begin{equation}
  c_{\text{fact}}(B) = \frac{\sum_{j: v_j = \text{Confirmed}} w_j \cdot c_j
    - \sum_{j: v_j = \text{Disputed}} w_j \cdot c_j}
    {\sum_{j} w_j}
\end{equation}

Claims with $c_{\text{fact}} > 0.8$ are classified as ``well-supported.''
Claims with $c_{\text{fact}} < -0.5$ are classified as ``disputed.''
Intermediate values indicate ongoing verification.

\section{Prediction Market Integration}

The markets\_integration zome connects knowledge claims to prediction
markets:

\begin{enumerate}[nosep]
  \item A prediction market is created for a specific claim (e.g.,
    ``Global average temperature will exceed 1.5C above pre-industrial
    by 2030'')
  \item Participants stake MYCEL on outcomes
  \item Market prices reflect collective confidence in the claim
  \item Resolved markets adjust participant reputation (correct predictions
    increase reputation, incorrect predictions decrease it)
\end{enumerate}

Market prices provide an independent signal for claim confidence,
complementing the expert verification pathway. The two signals are
combined with configurable weights:

\begin{equation}
  c_{\text{composite}} = \alpha \cdot c_{\text{fact}} +
    (1 - \alpha) \cdot c_{\text{market}}
\end{equation}

where $\alpha$ is the domain-specific weight for expert vs. market
confidence (default 0.7, prioritizing expert verification).

\section{Cross-Domain Corroboration}

\subsection{The Epistemic Mesh}

The Knowledge cluster's bridge zome enables cross-domain corroboration:

\begin{itemize}[nosep]
  \item A climate claim about deforestation can be corroborated by the
    Space cluster's satellite observations
  \item A supply chain provenance claim can be verified against the
    Energy cluster's grid measurements
  \item A health claim about nutrition can be checked against the
    Commons cluster's food production data
\end{itemize}

Cross-domain corroboration uses the \texttt{CrossDomainCorroboration}
evidence source type and adds Support/Contradict edges to the knowledge
graph.

\subsection{Cross-Cluster Query}

The query zome supports cross-cluster queries via the bridge:

\begin{lstlisting}
pub struct CrossClusterQuery {
    pub source_domain: String,
    pub target_cluster: CrossClusterRole,
    pub query: SemanticQuery,
    pub max_depth: u32,
    pub min_confidence: f64,
}
\end{lstlisting}

Queries are consciousness-gated: cross-cluster queries require Citizen+
tier to prevent abuse of the bridge dispatch mechanism.

\section{Consciousness-Weighted Epistemic Authority}

The key innovation of the epistemic mesh is that epistemic authority
derives from the consciousness profile rather than institutional position:

\begin{proposition}[Epistemic weight decomposition]
An agent's epistemic weight for a claim in domain $d$ is:
\begin{equation}
  w = \underbrace{E_d}_{\text{engagement}} \cdot
      \underbrace{f(\text{tier})}_{\text{governance}} \cdot
      \underbrace{R}_{\text{reputation}} \cdot
      \underbrace{(1 + \beta \cdot h_d)}_{\text{history}}
\end{equation}
where $E_d$ is domain engagement, $f(\text{tier})$ is the tier weight
function, $R$ is reputation, $h_d$ is the agent's verification accuracy
history in domain $d$, and $\beta$ is the history bonus factor.
\end{proposition}

This decomposition means that an agent's epistemic authority in a domain
is proportional to their actual engagement and track record, not their
credentials or institutional affiliation.

\section{Integration with Symthaea}

For agents connected to the Symthaea consciousness engine, the Knowledge
cluster provides a bidirectional integration:

\begin{itemize}[nosep]
  \item \textbf{Knowledge $\to$ Consciousness}: Verified knowledge claims
    feed the Knowledge Manager in Symthaea's cognitive loop, providing
    grounding for the reasoning engine.
  \item \textbf{Consciousness $\to$ Knowledge}: Symthaea's inference outputs
    (from the reasoning engine) can be registered as claims with appropriate
    provenance metadata, contributing to the collective knowledge graph.
\end{itemize}

\section{Evaluation}

The Knowledge cluster includes:
\begin{itemize}[nosep]
  \item Unit tests for claim validation, graph operations, and inference rules
  \item Integration tests for cross-zome query processing
  \item Property tests for confidence propagation bounds
  \item Cross-cluster integration tests for corroboration pathways
\end{itemize}

\section{Related Work}

\textbf{Google Knowledge Graph}~\cite{singhal2012}: Centralized, proprietary,
no verification mechanism.

\textbf{Wikidata}~\cite{vrandecic2014}: Open but centrally administered.
No confidence scoring or prediction market integration.

\textbf{The Graph}~\cite{thegraph2020}: Decentralized indexing for blockchain
data. Focuses on data availability, not knowledge verification.

\textbf{Augur}~\cite{peterson2015}: Decentralized prediction markets on
Ethereum. Addresses price discovery but not structured knowledge verification.

Our contribution is the integration of structured claims, logical inference,
expert verification, and prediction markets into a single epistemic mesh
with consciousness-weighted authority.

\section{Conclusion}

The epistemic mesh demonstrates that knowledge verification can be
decentralized without sacrificing rigor. Multiple independent verification
pathways---peer review, logical inference, market pricing, and cross-domain
corroboration---converge on confidence scores that reflect genuine epistemic
status rather than institutional authority. Consciousness-weighted verification
ensures that domain experts with demonstrated engagement contribute more to
the collective knowledge base, while the open participation model (Participant+
for claim registration) ensures that no gatekeeper can suppress inconvenient
truths.

\begin{thebibliography}{99}
  \bibitem{singhal2012} Singhal, A. (2012). Introducing the Knowledge Graph.
    Google Blog.
  \bibitem{vrandecic2014} Vrandecic, D. and Krotzsch, M. (2014). Wikidata:
    A Free Collaborative Knowledgebase. \textit{Communications of the ACM}.
  \bibitem{thegraph2020} The Graph (2020). The Graph: An Indexing Protocol
    for Querying Networks.
  \bibitem{peterson2015} Peterson, J. et al. (2015). Augur: A Decentralized
    Oracle and Prediction Market Platform. Whitepaper.
  \bibitem{brock2018} Brock, A. and Harris-Braun, E. (2018). Holochain:
    Scalable Agent-Centric Distributed Computing.
\end{thebibliography}

\end{document}
